\documentclass{article}
\usepackage[utf8]{inputenc}
\usepackage[spanish]{babel}
\usepackage{listings}
\usepackage{xcolor}
\usepackage{geometry}
\geometry{a4paper, margin=1in}

% Configuración para código C++
\definecolor{codegreen}{rgb}{0,0.6,0}        
\definecolor{codegray}{rgb}{0.5,0.5,0.5}     
\definecolor{codepurple}{rgb}{0.58,0,0.82}   
\definecolor{backcolour}{rgb}{0.95,0.95,0.92}

\lstset{
    backgroundcolor=\color{backcolour},      
    commentstyle=\color{codegreen},
    keywordstyle=\color{magenta},
    numberstyle=\tiny\color{codegray},       
    stringstyle=\color{codepurple},
    basicstyle=\ttfamily\small,
    breakatwhitespace=false,
    breaklines=true,
    captionpos=b,
    keepspaces=true,
    numbers=left,
    numbersep=5pt,
    showspaces=false,
    showstringspaces=false,
    showtabs=false,
    tabsize=2,
    language=C++,
    inputencoding=utf8,
    extendedchars=true,
    literate={á}{{\'a}}1 {é}{{\'e}}1 {í}{{\'i}}1 {ó}{{\'o}}1 {ú}{{\'u}}1 {ñ}{{\~n}}1 {¿}{{\textquestiondown}}1 {¡}{{\textexclamdown}}1
}

% Macros para las secciones de los apuntes
\newcommand{\quees}[1]{
    \section*{🤔 ¿QUÉ ES?}
    \hrulefill \\
    #1
}

\newcommand{\dichodeotraforma}[1]{
    \section*{💬 DICHO DE OTRA FORMA}
    \hrulefill \\
    #1
}

\newcommand{\diagrama}[1]{
    \section*{🎨 DIAGRAMA}
    \hrulefill \\
    #1
}

\newcommand{\comofunciona}[1]{
    \section*{⚙️ ¿CÓMO FUNCIONA?}
    \hrulefill \\
    #1
}

\newcommand{\complejidad}[1]{
    \section*{📱 COMPLEJIDAD (Big-O)}
    \hrulefill \\
    #1
}

\newcommand{\entrevistas}[1]{
    \section*{🎯 EN ENTREVISTAS}
    \hrulefill \\
    #1
}

\newcommand{\resumen}[1]{
    \section*{📋 RESUMEN RÁPIDO}
    \hrulefill \\
    #1
}

\begin{document}

\title{Linear Search (Búsqueda Lineal)}
\author{Aldo Zepeda Montoya}
\date{\today}
\maketitle

\subsection*{CATEGORÍA: 05_Algoritmos_Busqueda}

\quees{
Buscar tus llaves en casa 🔑 revisando cuarto por cuarto, cajón por
cajón, sin ningún orden especial. 

Empiezas en un lugar y revisas todo hasta que las encuentras o hasta
que verificas que no están en ningún lado. Es la forma más básica y 
natural de buscar algo.
}

\dichodeotraforma{
Linear Search es un algoritmo de búsqueda que recorre un conjunto de
datos secuencialmente (uno por uno) de principio a fin. 

No requiere que los datos estén ordenados. Su simplicidad lo hace
ideal para datasets pequeños o cuando solo vas a buscar algo una vez
y no vale la pena ordenar todo el conjunto primero.
}

\diagrama{\begin{verbatim}
Buscar el número 7 en el array:
[ 3, 1, 9, 7, 5 ]

Paso 1: ¿3 == 7? No.
Paso 2: ¿1 == 7? No.
Paso 3: ¿9 == 7? No.
Paso 4: ¿7 == 7? ¡SÍ! 🎉

Encontrado en 4 pasos.
\end{verbatim}}

\begin{lstlisting}[language=C++]
1. Empiezas en el primer elemento (índice 0).
2. Comparas el elemento actual con el valor que buscas.
3. Si son iguales, devuelves el índice o `true`.
4. Si no son iguales, pasas al siguiente elemento.
5. Repites hasta encontrarlo o llegar al final del array.
6. Si llegas al final y no lo encontraste, devuelves -1 o `false`.

📝 PSEUDOCÓDIGO
──────────────────────────────────────────────────────────────

función busquedaLineal(array, objetivo):
    PARA cada i desde 0 hasta tamaño(array) - 1:
        SI array[i] == objetivo:
            return i
    return -1

💻 CÓDIGO C++
──────────────────────────────────────────────────────────────

#include <iostream>
#include <vector>
using namespace std;

int linearSearch(const vector<int>& arr, int target) {
    for (int i = 0; i < arr.size(); i++) {
        if (arr[i] == target) {
            return i; // Encontrado, devolvemos el índice
        }
    }
    return -1; // No encontrado
}

int main() {
    vector<int> data = {10, 50, 30, 70, 80, 60, 20, 90, 40};
    int target = 20;

    int result = linearSearch(data, target);
    
    if (result != -1)
        cout << "Elemento encontrado en el índice: " << result << endl;
    else
        cout << "Elemento no encontrado." << endl;

    return 0;
}

⏱️ COMPLEJIDAD (Big-O)
──────────────────────────────────────────────────────────────

  Caso        │ Tiempo        │ Razón
  ────────────┼───────────────┼────────────────────────────────
  Best        │ O(1)          │ El elemento es el primero
  Average     │ O(n)          │ Tienes que revisar la mitad (n/2)
  Worst       │ O(n)          │ Está al final o no existe
  Espacio     │ O(1)          │ No usa memoria extra
\end{lstlisting}

\entrevistas{
✅ Cuándo usar: 
   - Datasets pequeños (n < 50).
   - Datos desordenados.
   - Búsquedas únicas (one-off).
   - Estructuras sin acceso aleatorio (como Linked Lists).

💡 Strings: Linear Search es la base para encontrar un caracter en 
   un string o una subcadena (algoritmo naive).

⚠️ Eficiencia: Si vas a hacer MUCHAS búsquedas sobre el mismo conjunto
   de datos, ¡no uses Linear Search! Mejor ordena los datos y usa 
   Binary Search, o guárdalos en un Hash Map.

🔑 Sentinel Search: Una pequeña optimización que elimina una de las
   dos comparaciones del loop (la que revisa si llegaste al final).
}

\resumen{
• Busca uno por uno.
• Muy simple de implementar.
• No requiere orden previo.
• Tiempo lineal O(n).
• Memoria constante O(1).
• Ineficiente para datasets grandes.
════════════════════════════════════════════════════════════════
}

\end{document}